% !TEX program = xelatex
% ============================================================================
% Manual de instalacion y ejecucion (PostgreSQL local, pgAdmin, Neon y Backend)
% Proyecto Final de Bases de Datos - Inventario, Compras y Ventas
% ============================================================================
\documentclass[11pt,oneside]{article}

\usepackage{estilo-proyecto}
\docnombre{Manual de instalaci\'on y ejecuci\'on}

\addto\captionsspanish{\renewcommand{\contentsname}{Contenido}}
\setcounter{tocdepth}{2}

\begin{document}

\portadilla{Proyecto final \textperiodcentered{} Bases de datos \textperiodcentered{} Manual t\'ecnico}
  {Manual de instalaci\'on y ejecuci\'on}
  {Despliegue local en PostgreSQL (psql / pgAdmin), Neon y microservicios Spring Boot}

\tableofcontents
\vspace{0.5em}

% ============================================================================
\section{Requisitos previos}

\begin{itemize}[leftmargin=1.5em, itemsep=3pt]
  \item \textbf{PostgreSQL 14 o superior} (validado en \textbf{PostgreSQL 18.6})
        con los binarios \texttt{initdb}, \texttt{pg\_ctl} y \texttt{psql}, m\'as
        la extensi\'on est\'andar \texttt{postgres\_fdw}.
  \item \textbf{pgAdmin 4} (opcional) para ejecuci\'on gr\'afica de los
        scripts SQL y consulta visual de tablas, particiones y vistas.
  \item \textbf{JDK 25} (opcional, solo para compilar y ejecutar los
        microservicios de \texttt{Backend/}); Maven Wrapper (\texttt{./mvnw})
        ya viene incluido.
\end{itemize}

% ============================================================================
\section{Organizaci\'on de los scripts SQL}

En la carpeta \texttt{Scripts sql/} cada microservicio tiene su propio
subdirectorio con cuatro scripts numerados en orden estricto de ejecuci\'on:

\begin{center}
\renewcommand{\arraystretch}{1.22}
{\small
\begin{tabularx}{\linewidth}{@{} >{\ttfamily}p{4.8cm} >{\raggedright\arraybackslash}X @{}}
\toprule
\rowcolor{fondo}
{\sffamily\bfseries Subdirectorio / Archivo} & {\sffamily\bfseries Funci\'on} \\
\midrule
catalogo\_terceros/          & 01\_tablas, 02\_constraints, 03\_datos, 04\_seguridad \\
compras\_abastecimiento/     & 01\_tablas, 02\_constraints, 03\_datos, 04\_seguridad \\
inventarios\_transferencias/ & 01\_tablas (partici\'on LIST), 02\_constraints, 03\_datos, 04\_seguridad \\
ventas\_pagos/               & 01\_tablas (partici\'on LIST), 02\_constraints, 03\_datos, 04\_seguridad \\
identidad\_accesos/          & 01\_tablas, 02\_constraints, 03\_datos, 04\_seguridad \\
reportes/                    & 01\_fdw\_integracion, 02\_vistas, 03\_consultas\_analiticas, 04\_seguridad, ejecutar\_pruebas.sh \\
local-db.sh                  & Script automatizado para cl\'uster local en puerto 5433 \\
\bottomrule
\end{tabularx}}
\end{center}

% ============================================================================
\section{Ejecuci\'on r\'apida en local con \texttt{local-db.sh}}

El proyecto incluye \texttt{Scripts sql/local-db.sh}, que crea un cl\'uster
PostgreSQL aislado en \texttt{.pgdata/} (puerto \texttt{5433}, usuario
\texttt{postgres}) para no interferir con otras instalaciones del equipo:

\begin{verbatim}
cd "Scripts sql"
./local-db.sh init        # Arranca PostgreSQL (5433), crea las 6 BD y carga todo
./local-db.sh pruebas     # Ejecuta la matriz de 10 pruebas de verificación
./local-db.sh consultas   # Ejecuta las 20 consultas analíticas en reportes
./local-db.sh status      # Consulta estado y lista de bases de datos
./local-db.sh stop        # Detiene el clúster local
\end{verbatim}

% ============================================================================
\section{Instalaci\'on manual paso a paso (\texttt{psql} o pgAdmin)}

Si prefiere usar una instancia existente de PostgreSQL o \textbf{pgAdmin 4},
siga este orden exacto:

\begin{enumerate}[leftmargin=1.5em, itemsep=4pt]
  \item \textbf{Crear las seis bases de datos:}
\begin{verbatim}
CREATE DATABASE catalogo_terceros ENCODING 'UTF8';
CREATE DATABASE compras_abastecimiento ENCODING 'UTF8';
CREATE DATABASE inventarios_transferencias ENCODING 'UTF8';
CREATE DATABASE ventas_pagos ENCODING 'UTF8';
CREATE DATABASE identidad_accesos ENCODING 'UTF8';
CREATE DATABASE reportes ENCODING 'UTF8';
\end{verbatim}

  \item \textbf{Cargar cada base operacional:} con\'ectese a cada una de las
        cinco primeras bases (en \texttt{psql -d <base>} o mediante \textit{Query
        Tool} en pgAdmin) y ejecute en orden:
        \texttt{01\_tablas.sql} $\to$ \texttt{02\_constraints.sql} $\to$
        \texttt{03\_datos.sql} $\to$ \texttt{04\_seguridad.sql}.
        \begin{verbatim}
for db in catalogo_terceros compras_abastecimiento \
          inventarios_transferencias ventas_pagos identidad_accesos; do
  for f in 01_tablas.sql 02_constraints.sql 03_datos.sql 04_seguridad.sql; do
    psql -h localhost -p 5433 -U postgres -d "$db" -v ON_ERROR_STOP=1 -f "$db/$f"
  done
done
        \end{verbatim}

  \item \textbf{Configurar \texttt{reportes}:} con\'ectese a la base
        \texttt{reportes} y ejecute:
        \begin{enumerate}[label=\alph*)]
          \item \texttt{reportes/01\_fdw\_integracion.sql} (crea servidores
                \texttt{postgres\_fdw} e importa tablas y particiones).
          \item \texttt{reportes/02\_vistas.sql} (crea las 10 vistas
                obligatorias m\'as 2 vistas \texttt{UNION ALL} de fragmentos:
                12 vistas en total).
          \item \texttt{reportes/04\_seguridad.sql} (crea roles y permisos de
                solo lectura, incluida la cuenta de servicio \texttt{usr\_reportes\_app}).
          \item \texttt{reportes/03\_consultas\_analiticas.sql} (ejecuta las
                20 consultas numeradas).
        \end{enumerate}
\end{enumerate}

% ============================================================================
\section{Despliegue en la nube (Neon + Render)}

\begin{itemize}[leftmargin=1.5em, itemsep=3pt]
  \item \textbf{Bases de datos en Neon:} cree un proyecto o base por cada
        microservicio y ejecute sus scripts \texttt{01}--\texttt{04}. En
        \texttt{reportes/01\_fdw\_integracion.sql}, sustituya \texttt{host},
        \texttt{port} (\texttt{5432}), \texttt{dbname} y credenciales de
        \texttt{CREATE SERVER} / \texttt{USER MAPPING} por los de Neon con
        \texttt{sslmode 'require'} (y omita \texttt{password\_required 'false'},
        ya que \texttt{neondb\_owner} no es superusuario). Como alternativa
        automatizada, puede usar \texttt{Scripts sql/neon-db.sh} leyendo
        \texttt{claves\_neon.txt} (\texttt{init}, \texttt{pruebas},
        \texttt{consultas}, \texttt{env}, \texttt{status}).
  \item \textbf{Microservicios en Render:} el archivo \texttt{render.yaml} en
        la ra\'iz define los 6 servicios web (\textit{New $\to$ Blueprint}) y
        solicita \texttt{DB\_URL}, \texttt{DB\_USER} y \texttt{DB\_PASSWORD} de
        cada base.
\end{itemize}

% ============================================================================
\section{Ejecuci\'on local del backend Spring Boot}

\begin{enumerate}[leftmargin=1.5em, itemsep=3pt]
  \item Copie \texttt{Backend/.env.example} a \texttt{Backend/.env} y ajuste
        las cadenas JDBC al cl\'uster local (\texttt{localhost:5433}) o a Neon.
        Para \texttt{reportes-service}, use la cuenta de solo lectura
        \texttt{usr\_reportes\_app} (\texttt{Reportes\#App\_2026}).
  \item Levante los seis servicios (puertos \texttt{8081}--\texttt{8086}); en
        arranque en fr\'io \texttt{run-all.sh} espera autom\'aticamente la
        apertura de \texttt{/api/health} (hasta $\sim$80\,s):
\begin{verbatim}
cd Backend
./run-all.sh start --build    # Compila JARs y arranca 8081..8086
./run-all.sh status           # Verifica estado de procesos
./run-all.sh stop             # Detiene los 6 microservicios
\end{verbatim}
  \item Verifique la integraci\'on consultando:
\begin{verbatim}
curl http://localhost:8086/api/health
curl http://localhost:8086/api/reportes/vistas
curl http://localhost:8086/api/reportes/corporativo
curl http://localhost:8086/api/reportes/alertas
\end{verbatim}
\end{enumerate}

\end{document}
